% Generated draft: requires mathematical and source review.
\documentclass[11pt]{article}
\usepackage[margin=25mm]{geometry}
\usepackage{fontspec}
\setmainfont{Latin Modern Roman}
\usepackage{amsmath,amssymb,mathtools,unicode-math}
\setmathfont{Latin Modern Math}
\usepackage{enumitem,needspace,xurl,hyperref}
\setlength{\parindent}{0pt}
\setlength{\parskip}{5pt}
\setlength{\emergencystretch}{3em}
\newcommand{\N}{\mathbb{N}}
\newcommand{\Z}{\mathbb{Z}}
\newcommand{\Q}{\mathbb{Q}}
\newcommand{\R}{\mathbb{R}}
\newcommand{\C}{\mathbb{C}}
\newcommand{\F}{\mathbb{F}}
\newcommand{\A}{\mathbb{A}}
\newcommand{\PP}{\mathbb{P}}
\newcommand{\E}{\mathbb{E}}
\newcommand{\eps}{\varepsilon}
\newcommand{\dd}{\,\mathrm{d}}
\begin{document}
\section*{UnsolvedMath ID: 30003402}
\textbf{Source problem code:} catalogue entry
\subsection{Definitions and mathematical statement}
A linear order is a set equipped with a total order. Every subset is given the induced order. A map \(f:L\to K\) is order-preserving if \(a\le_L b\) implies \(f(a)\le_K f(b)\). For a nonempty linear order \(L\), write \(\operatorname{SS}(L)\) if, for every nonempty subset \(K\subseteq L\), there is an order-preserving surjection \(f:L\twoheadrightarrow K\). This is the strong-surjectivity property used below. The empty-order convention is not supplied and is discussed explicitly below.

Write \(\mathrm{MA}_{\aleph_1}\) for Martin's axiom for families of at most \(\aleph_1\) dense subsets of a partial order satisfying the countable chain condition. Write \(\mathrm{PFA}\) for the proper forcing axiom: every family of at most \(\aleph_1\) dense subsets of a proper partial order is met by a filter. A Countryman line is a linear order \(C\) of cardinality \(\aleph_1\) such that \(C\times C\), ordered coordinatewise, is the union of countably many chains. These are standard terminology conventions; the supplied record does not itself give these definitions.

The record asks the following questions. For the consistency questions, \(\mathrm{ZFC}\) is the conventional background theory adopted here; the record does not explicitly specify a base theory or relative-consistency strength.
\begin{enumerate}
\item Is the theory
\[
\mathrm{ZFC}+\exists L\,\bigl(\operatorname{SS}(L)\land |L|>\aleph_2\bigr)
\]
consistent, where \(L\) ranges over nonempty linear orders?

\item Is it consistent that a strongly surjective linear order contains no minimal strongly surjective suborder? The word ``minimal'' is not defined in the supplied record. In particular, the record does not specify the comparison relation, whether suborders must be uncountable, or whether minimality is taken up to order isomorphism. Consequently this question cannot be given a unique precise quantified formulation from the supplied evidence. If singleton suborders are allowed and minimality means minimality under inclusion among nonempty strongly surjective suborders, every nonempty order contains such a minimal suborder. This observation identifies a missing qualification; it does not identify the intended formulation.

\item Does \(\mathrm{ZFC}+\mathrm{MA}_{\aleph_1}\) imply
\[
\exists L\subseteq\R\,\bigl(|L|>\aleph_0\land\operatorname{SS}(L)\bigr),
\]
where \(L\) has the usual real order restricted to \(L\)?

\item For every nonempty strongly surjective linear order \(L\) and every \(x\in L\), must \(L\setminus\{x\}\), with its induced order, be strongly surjective? With the nonempty-order definition above, the substantive assertion is
\[
\forall L\,\left[\operatorname{SS}(L)\Longrightarrow
\forall x\in L\,\left(L\setminus\{x\}\ne\varnothing
\Longrightarrow\operatorname{SS}(L\setminus\{x\})\right)\right].
\]
The source also quantifies over deletion from singleton orders, but does not specify whether the empty order is admitted or called strongly surjective. That boundary case remains unresolved in this draft.

\item Separately for each \(A\in\{\mathrm{MA}_{\aleph_1},\mathrm{PFA}\}\), does \(\mathrm{ZFC}+A\) imply
\[
\exists C\,\bigl(C\text{ is a Countryman line}\land\operatorname{SS}(C)\bigr)?
\]
The wording ``exist under'' is interpreted here as an implication from the stated axiom; the record does not explicitly distinguish this reading from mere consistency together with that axiom.
\end{enumerate}

The supplied registry labels the collection of questions as partially solved. Its dated assessment mentions strongly surjective Countryman lines under Martin's axiom and says that size, minimality, deletion, and other consistency questions remain open. These are supplied curation claims, not independently verified mathematical conclusions. In particular, the assessment does not specify which version of Martin's axiom supports the cited Countryman-line result or give an item-by-item resolution.
\subsection{Short English statement}
A strongly surjective order maps onto each of its nonempty suborders by an order-preserving surjection. The questions concern possible cardinalities, minimal suborders, uncountable examples inside \(\R\), deletion of one point, and Countryman examples under forcing axioms. Minimality, the empty-order boundary case, and the exact scope of the supplied partial-resolution assessment require review.
\subsection{Sources}
\begin{itemize}
\item Local UnsolvedMath registry, supplied complete record: permanent ID \(30003402\), source problem code \texttt{catalogue entry}, title ``Consistency and Structure of Strongly Surjective Linear Orders,'' in the Oberwolfach Reports Open Problems collection. No individual proposer is supplied; the proposed year is recorded as 2017.
\item Supplied source citation: \emph{Set Theory (2018)}, Oberwolfach Reports, DOI 10.4171/owr/2017/11, pp. 527--589. \url{https://doi.org/10.4171/owr/2017/11}. The title's year, the DOI's year component, and the recorded proposed year are retained without reconciliation.
\item Literature source labeled ``Oberwolfach report'': \url{https://ems.press/content/serial-article-files/46674}.
\item Literature source labeled ``Soukup'': \url{https://doi.org/10.1007/s11083-018-9454-7}. No full bibliographic entry or precise supporting theorem is supplied.
\item Literature source labeled ``Polymeris--Martinez-Ranero'': \url{https://arxiv.org/abs/2503.13728}. No full bibliographic entry or precise supporting theorem is supplied.
\item Literature assessment date, copied exactly: \texttt{2026-08-21}. The supplied dated curation labels the problem partially solved, mentions strongly surjective Countryman lines under Martin's axiom, and describes several remaining questions as open. This assessment is supplied dated curation, not independently verified. The background explicitly requests independent verification of time-sensitive status and cited claims.
\item The registry's supplied statement-status label is \texttt{corrected\_verified}; its accompanying explanation concerns preservation of mathematical content during repairs to notation, typography, and sentence structure. This label does not establish independent mathematical verification of this draft.
\end{itemize}
\end{document}
